% ============================================================
% Chapter 6 -- Learning the parameters: EM on the chain
% ============================================================

\chapter{Learning the Parameters: EM on the Chain}\label{ch:em-method}

Everything so far assumed the transition kernel was given. This chapter drops
that assumption: the kernel is unknown and is recovered from the noised
sequences alone, with no clean data, no supervision, and no differentiation
through the inference. The tool is expectation--maximisation, and the reason
it is cheap here is Fisher's identity: at each parameter value one
forward--backward pass of belief propagation already contains the exact
gradient of the marginal log-likelihood, so every E-step is the recursion of
Chapter~\ref{ch:gaussian-bp} with the kernel held fixed, and nothing is ever
back-propagated through it.

\section{Model: an estimated mixture innovation}\label{sec:em-model}

The clean chain is \eqref{eq:l-chain-prior} in the index convention of this
chapter, $a_1 \sim \Nd(0,1)$ and $a_i = \alpha a_{i-1} + \varepsilon_i$ for
$i = 2,\dots,\nsites$, with i.i.d.\ centred innovations of variance
$\ivar = 1-\alpha^2$ now drawn from a member of a parametric family; noising
is the channel \eqref{eq:ou-channel}. The estimated family is a
$C$-component Gaussian mixture on the innovation,
\begin{equation}
  \kernel_\theta(a' \mid a) = \sum_{c=1}^{C} \pi_c \,
    \Nd\!\big(a' - \alpha a;\ \nu_c,\ s_c^2\big),
  \qquad \theta = \{\alpha, (\pi_c,\nu_c,s_c)_{c=1}^C\},
  \label{eq:em-mixture-kernel}
\end{equation}
with $\pi_c \ge 0$, $\sum_c\pi_c = 1$, $s_c^2 \ge s_{\min}^2$, $\nu_c \in \R$
and $\alpha \in (-1,1)$, so that both the autoregressive coefficient and the
innovation shape are estimated jointly by maximum marginal likelihood from
noised data; the count of free parameters is $3C$, and $C = 1$ is the Gaussian
chain of Chapter~\ref{ch:gaussian}. The fitted family deliberately enforces
neither $\sum_c\pi_c\nu_c = 0$ nor a matched total variance, although the
generating innovations satisfy both: projecting the means back to zero after
each update broke monotone ascent by $\sim 10^{-2}$ nats per edge when tried,
whereas the unconstrained update is exactly monotone and the fitted innovation
mean lands at the $10^{-3}$ level on its own, so the moment conditions are
reported as diagnostics rather than imposed. The constraint $|\alpha| < 1$
never binds in any fit reported here, and component labels are exchangeable,
so label-dependent diagnostics use a fixed ordering.

\section{The EM algorithm on the chain}\label{sec:em-algorithm}

Observed is $x = (x_1,\dots,x_\nsites)$; latent are the clean trajectory $a$
and, for every transition, the label $c_i \in \{1,\dots,C\}$ of the component
that generated its innovation. Maximum likelihood targets
$\mathcal L(\theta) = \log p_\theta(x) = \log\sum_c\int p_\theta(a,c,x)\,\dd a$,
a logarithm outside an integral over the trajectory and a sum over labels.
With $q^{(k)}(a,c) := p_{\theta^{(k)}}(a,c\mid x)$ the posterior at the
current iterate and the surrogate
$Q(\theta\mid\theta^{(k)}) := \E_{q^{(k)}}[\log p_\theta(a,c,x)]$, taking
$\E_{q^{(k)}}$ of $\log p_\theta(x) = \log p_\theta(a,c,x) - \log p_\theta(a,c\mid x)$
gives the exact decomposition
\begin{equation}
  \mathcal L(\theta)
  \;=\;
  Q(\theta\mid\theta^{(k)}) + \mathcal H(q^{(k)})
  + \operatorname{KL}\!\big( q^{(k)} \,\big\|\, p_\theta(a,c\mid x) \big),
  \label{eq:em-decomposition}
\end{equation}
a lower bound $Q + \mathcal H$ that touches $\mathcal L$ at $\theta^{(k)}$,
so any $\theta^{(k+1)}$ that increases $Q$ increases $\mathcal L$. The E-step
computes $q^{(k)}$; the M-step increases $Q$. Because the posterior factor
graph is a tree, one forward and one backward sweep return the exact
single-site and pairwise posterior marginals at the level of functional
messages: the E-step needs neither a variational approximation nor Monte
Carlo, and is exact under the model before the grid is introduced.

Conditional on $c_i = c$ the transition is a single Gaussian,
$\widetilde\kernel_\theta(a',c\mid a) = \pi_c\Nd(a'-\alpha a;\nu_c,s_c^2)$, and
only the transition factors of the complete joint law
$p(a_1)\prod_i\widetilde\kernel_\theta(a_i,c_i\mid a_{i-1})\prod_i p(x_i\mid a_i)$
depend on $\theta$:
\begin{equation}
  \log p_\theta(a,c,x)
  \;\overset{c}{=}\;
  \sum_{i=2}^{\nsites}\sum_{c=1}^{C}\mathbf 1\{c_i=c\}
  \Big[ \log\pi_c - \tfrac12\log s_c^2
        - \frac{(a_i-\alpha a_{i-1}-\nu_c)^2}{2s_c^2} \Big],
  \label{eq:em-complete-ll}
\end{equation}
$\overset{c}{=}$ denoting equality up to $\theta$-independent terms. Exposing
the label turns the log of a mixture into a weighted sum of Gaussian
log-densities with closed-form maximisers. Its expectation needs, for each
edge, only the pairwise posterior supplied by the two sweeps,
\begin{equation}
  b_i(a_{i-1},a_i) \;\propto\;
  \lmsg_{i-1}(a_{i-1})\,\phi_{i-1}(a_{i-1})\,
  \kernel_{\theta^{(k)}}(a_i\mid a_{i-1})\,
  \phi_i(a_i)\,\rmsg_i(a_i),
  \qquad \phi_i(a_i) = p(x_i \mid a_i),
  \label{eq:em-pairwise}
\end{equation}
and the responsibility $r_{i,c} = \Prob_{\theta^{(k)}}(c_i = c \mid a_{i-1}, a_i)$
of component $c$ given the pair, Bayes' rule on the finite mixture at the
current parameters, which needs no further recursion. Writing
$\langle\cdot\rangle_i$ for expectation under $b_i$, the E-step reduces to six
sufficient statistics per component, summed over all edges of all sequences,
\begin{equation}
  \begin{aligned}
    R_c &= \textstyle\sum_i \langle r_{i,c}\rangle_i, \qquad
    U_c = \textstyle\sum_i \langle r_{i,c}\,a_i\rangle_i, \qquad
    V_c = \textstyle\sum_i \langle r_{i,c}\,a_{i-1}\rangle_i, \\
    P_c &= \textstyle\sum_i \langle r_{i,c}\,a_{i-1}a_i\rangle_i, \qquad
    W_c = \textstyle\sum_i \langle r_{i,c}\,a_{i-1}^2\rangle_i, \qquad
    G_c = \textstyle\sum_i \langle r_{i,c}\,a_i^2\rangle_i,
  \end{aligned}
  \label{eq:em-suffstats}
\end{equation}
in terms of which, with
$S_c(\alpha,\nu_c) := G_c - 2\alpha P_c + \alpha^2 W_c - 2\nu_c U_c + 2\alpha\nu_c V_c + \nu_c^2 R_c$,
\begin{equation}
  Q(\theta\mid\theta^{(k)})
  \;\overset{c}{=}\;
  \sum_{c=1}^{C}
  \Big[ R_c\log\pi_c - \frac{R_c}{2}\log s_c^2 - \frac{S_c(\alpha,\nu_c)}{2s_c^2} \Big].
  \label{eq:em-Q-suffstats}
\end{equation}

\paragraph{Fisher's identity.} Differentiating $\log p_\theta(x)$ under the
integral, licit because the Gaussian factors and the mixture kernel make
$\nabla_\theta p_\theta(a,c,x)$ dominated on the compact parameter set every
fit stays in,
\begin{equation}
  \nabla_\theta\mathcal L(\theta)
  = \E_{a,c\mid x,\theta}\big[\nabla_\theta\log p_\theta(a,c,x)\big]
  = \E_{a,c\mid x,\theta}\Big[\sum_{i=2}^{\nsites}\nabla_\theta\log
      \widetilde\kernel_\theta(a_i,c_i\mid a_{i-1})\Big],
  \label{eq:em-fisher}
\end{equation}
so that $\nabla_\theta\mathcal L(\theta)|_{\theta^{(k)}} = \nabla_\theta
Q(\theta\mid\theta^{(k)})|_{\theta^{(k)}}$.
BP is used only to compute the posterior; once computed it is held fixed, and
the derivative acts on the candidate kernel inside $\log\widetilde\kernel_\theta$,
not on the messages that produced the posterior. BP supplies posterior
expectations, Fisher's identity converts them into the likelihood gradient,
and no derivative is taken through the inference.

\paragraph{M-step.} Setting the derivatives of \eqref{eq:em-Q-suffstats} to
zero block by block, with a multiplier for $\sum_c\pi_c = 1$,
\begin{equation}
  \boxed{\;
  \pi_c^{+} = \frac{R_c}{\sum_d R_d},
  \qquad
  \nu_c^{+} = \frac{U_c-\alpha V_c}{R_c},
  \qquad
  (s_c^2)^{+} = \frac{G_c-2\alpha P_c+\alpha^2W_c - R_c(\nu_c^{+})^2}{R_c},
  \;}
  \label{eq:em-mstep-mixture}
\end{equation}
and, with the mixture block fixed,
\begin{equation}
  \boxed{\;
  \alpha^{+}
  = \frac{\sum_c s_c^{-2}\big(P_c-\nu_cV_c\big)}{\sum_c s_c^{-2}W_c},
  \;}
  \label{eq:em-mstep-alpha}
\end{equation}
a precision-weighted regression in which components with smaller innovation
variance carry more information about $\alpha$; at $C = 1$, $\nu_1 = 0$ it is
the posterior analogue of the least-squares slope,
$\sum_i\E[a_{i-1}a_i\mid x]/\sum_i\E[a_{i-1}^2\mid x]$. The iteration is
\begin{equation}
  \boxed{\;
  \text{E-step:}\ \theta^{(k)} \xrightarrow{\ \mathrm{BP}\ }
  \{R_c,U_c,V_c,P_c,W_c,G_c\}_{c=1}^{C}
  \qquad
  \text{M-step:}\ \{R_c,\dots,G_c\} \xrightarrow{\ \eqref{eq:em-mstep-mixture},\,\eqref{eq:em-mstep-alpha}\ }
  \theta^{(k+1)} .
  \;}
  \label{eq:em-summary}
\end{equation}
The updates are conditional maximisers, so the implementation is an ECM
scheme \citep{meng1993maximum}; by \eqref{eq:em-decomposition} any increase
of $Q$ increases $\mathcal L$, and the two blocks are alternated until the
gain in $Q$ per edge falls below tolerance or sixteen inner sweeps are
reached. The cap is chosen from the numerics: at
$C = 8$, $\nseq = 256$ and $40$ outer iterations, raising it from four to
sixteen moves the fitted excess kurtosis from $1.68$ to $2.14$ while
$\widehat\alpha$ changes in the fourth decimal, so the shape converges far
more slowly than the correlation, and by roughly $60$ outer iterations the
inner budget has little remaining effect. Two guards keep
\eqref{eq:em-mstep-mixture} defined: the responsibility mass is floored at
$10^{-300}$ before division and the variance at $s_c^2 \ge 10^{-6}$, far
below the grid spacing, so the floor marks degeneracy rather than shaping a
resolved fit; a separate reporting gate (\provenanceappendix) excludes fits
whose narrowest component is under two grid cells wide. The derivation
concerns the continuum model; the implementation evaluates densities on the
grid of Section~\ref{sec:em-numerical}, so the objective whose monotone
increase Table~\ref{tab:em-checks} certifies is the quadrature-defined
surrogate of the discretised factor model.

\section{The numerical representation}
\label{sec:em-numerical}

As in Section~\ref{sec:l-grid}, messages live on $\ngrid$ equispaced points of
$[-A,A]$ with trapezoidal integrals; the pairwise belief is an
$\ngrid\times\ngrid$ array, and one EM iteration costs $O(\nsites\ngrid^2)$ per
sequence with a factor $C$ only in the responsibilities. Two numerical errors
are kept apart (Table~\ref{tab:em-grid}, Figure~\ref{fig:em-grid}).
\emph{Truncation}, from replacing $\R$ by $[-A,A]$, is monitored through the
forward-message boundary mass and, because it depends on the sampled
trajectory, reported over $\gdchains$ chains, whose batch maximum exceeds a
single-chain reading by a median factor $\gdboundmed$ and by as much as
$\gdboundhi$. \emph{Quadrature}, from the spacing $h = 2A/(\ngrid-1)$, is
$O(h^2)$ in the interior, the residual falling by a factor $4.00$ per halving,
and does not converge near the boundary because the remaining error there is
truncation. The transition array is the continuum density evaluated on the
grid and integrated with trapezoidal weights \emph{without} renormalising each
parent-conditioned column: BP is then exact for the resulting discretised
factor model, which is not exactly a normalised finite-state kernel, and the
discrepancy is the column-mass residual of Table~\ref{tab:em-grid}. Column
renormalisation would introduce a $\theta$-dependent normaliser $Z_\theta$
whose gradient would enter Fisher's identity and destroy the closed-form
M-step, so the convention fits the continuum kernel through the quadrature
rule, with truncation controlled separately. Posterior-mean errors are
amplified in the score by $e^{-t}/\Deltat \sim 1/2t$, so numerical and
estimation deficiencies are most visible at small $t$.

\begin{center}
\begin{minipage}{\linewidth}\centering
\captionof{table}[The two discretisation budgets]{Numerical error
diagnostics. \emph{Upper:} boundary mass against half-width at $\ngrid=401$,
over $\gdchains$ chains. \emph{Lower:} column-mass residual against spacing
at $A=8$. Working choice: $\ngrid=401$, $A=8$.}
\label{tab:em-grid}
\small
\begin{tabular}{@{}lcccc@{}}
\toprule
\multicolumn{5}{@{}l}{\emph{truncation: boundary mass against half-width} ($\ngrid=401$)}\\
half-width $A$      & $4$ & $6$ & $8$ & $10$ \\
\midrule
max over chains     & $2.3\times10^{-3}$ & $1.6\times10^{-4}$ & $1.2\times10^{-6}$ & $3.2\times10^{-8}$ \\
$90$th percentile   & $2.7\times10^{-4}$ & $2.5\times10^{-6}$ & $1.5\times10^{-8}$ & $7.8\times10^{-11}$ \\
median              & $3.8\times10^{-5}$ & $3.7\times10^{-7}$ & $3.0\times10^{-9}$ & $1.8\times10^{-11}$ \\
max / median        & $61\times$ & $434\times$ & $392\times$ & $1796\times$ \\
\midrule
\multicolumn{5}{@{}l}{\emph{quadrature: column-mass residual against spacing} ($A=8$)}\\
$\ngrid$ \ / \ $h$  & \multicolumn{1}{c}{$201$ / $0.080$} & \multicolumn{1}{c}{$401$ / $0.040$} & \multicolumn{1}{c}{$801$ / $0.020$} & rate \\
\midrule
interior column     & $3.84\times10^{-3}$ & $9.61\times10^{-4}$ & $2.40\times10^{-4}$ & $4.00\times$ per halving \\
edge column         & $1.62\times10^{-2}$ & $1.90\times10^{-2}$ & $1.97\times10^{-2}$ & none \\
\bottomrule
\end{tabular}
\end{minipage}
\end{center}

\begin{figure}[t]\centering
\includegraphics[width=\linewidth]{fig_em_grid.pdf}
\caption[The two discretisation budgets, separated]{\emph{(a)} Boundary mass
falls rapidly with the half-width $A$ but varies substantially across
chains. \emph{(b)} Interior column-mass error follows the $O(h^2)$ rate; the
edge error is truncation-dominated. \emph{(c)} Relative score error against
diffusion time for three resolutions; failure sets in where
\eqref{eq:l-resolution} is violated.}
\label{fig:em-grid}
\end{figure}

\section{Diagnostics, and what EM recovers}\label{sec:em-diagnostics}

The recursion agrees with independent references to machine precision and the
surrogate never decreases in any logged run (Table~\ref{tab:em-checks});
$\widehat\alpha$ converges to the same value from initialisations spanning
$[-0.4, 0.6]$, so it is estimated rather than supplied, and its RMSE and that
of the innovation variance decay approximately at the $\nseq^{-1/2}$ rate
(Figure~\ref{fig:em-diagnostics}).

\begin{center}
\begin{minipage}{\linewidth}\centering
\captionof{table}[Correctness checks on the EM implementation]{Correctness
and monotonicity checks at $\ngrid=401$, $A=8$.}
\label{tab:em-checks}
\small
\setlength{\tabcolsep}{4pt}
\begin{tabular}{@{}llc@{}}
\toprule
check & reference & discrepancy \\
\midrule
posterior mean & closed-form Gaussian chain (Chapter~\ref{ch:gaussian})
  & $9.2\times10^{-15}$ \\
posterior mean & brute-force enumeration, independent code
  & $1.6\times10^{-14}$ \\
log-evidence   & brute-force enumeration, independent code
  & $1.8\times10^{-15}$ \\
\midrule
min.\ increase of $Q$, twelve steps & Gaussian kernel & $+2.61$ \\
 & Laplace kernel & $+4.34$ \\
 & $C=8$ mixture kernel & $+4.18$ \\
decreases in any logged run & & $0$ \\
\bottomrule
\end{tabular}
\end{minipage}
\end{center}

\begin{figure}[t]\centering
\includegraphics[width=\linewidth]{fig_em_diagnostics.pdf}
\caption[Optimisation diagnostics for the EM procedure]{Diagnostics at
$C=4$. \emph{(a)} Marginal log-likelihood increases monotonically from four
initialisations of $\alpha$. \emph{(b)} $\widehat\alpha$ converges to
$[0.8517,0.8520]$ from starts spanning $[-0.4,0.6]$, against the true
$0.85$. \emph{(c)} RMSE of $\widehat\alpha$ and of the innovation variance
against $\nseq$ over $16$ seeds, with an $\nseq^{-1/2}$ reference.}
\label{fig:em-diagnostics}
\end{figure}

The parameters are not recovered equally well (Table~\ref{tab:em-recovery},
Figure~\ref{fig:em-innovation}): at $C=8$ and $\nseq\ge512$ the correlation
is recovered to below $0.1\%$ relative error, the innovation variance to the
percent level, and the fourth cumulant only to tens of percent. Each row is a
single fit and the shape column is not monotone in $\nseq$, the seed-to-seed
spread of the fourth cumulant being of the order of its error, so the rows
demonstrate the hierarchy, not a convergence law for the shape. The ordering
is what the channel predicts: by \eqref{eq:l-state-kurtosis} the fourth
cumulant is attenuated as $e^{-4t}$ while the second moment is preserved, and
a finite Gaussian mixture also rounds the Laplace cusp, so the residual shape
error mixes information lost through the channel with approximation by the
family. ``Recovery'' means numerical agreement over the studied regime; no
identifiability or consistency theorem is claimed for the mixture family. The
estimated kernel is therefore most reliable in its correlation structure,
reasonably accurate in scale, and substantially less informative about tail
shape.

\begin{center}
\begin{minipage}{\linewidth}\centering
\captionof{table}[What EM recovers from noised sequences]{Estimates after
$61$ EM iterations on noised Laplace chains ($\alpha=0.85$,
$\ivar=0.2775$, excess kurtosis $3.0$, $\nsites=32$).}
\label{tab:em-recovery}
\small
\setlength{\tabcolsep}{4pt}
\begin{tabular}{@{}ccccc ccc@{}}
\toprule
$C$ & $\nseq$ & $\widehat\alpha$ & $\widehat\ivar$ & $\widehat\kappa_4$
& $|\Delta\alpha|/\alpha$ & $|\Delta\ivar|/\ivar$ & $|\Delta\kappa_4|/\kappa_4$ \\
\midrule
$4$ & $128$  & $0.8649$ & $0.2686$ & $1.61$ & $1.8\%$  & $3.2\%$ & $46\%$ \\
$4$ & $512$  & $0.8518$ & $0.2844$ & $2.27$ & $0.2\%$  & $2.5\%$ & $24\%$ \\
$4$ & $2048$ & $0.8465$ & $0.2818$ & $1.82$ & $0.4\%$  & $1.5\%$ & $39\%$ \\
\addlinespace
$8$ & $128$  & $0.8294$ & $0.2824$ & $3.18$ & $2.4\%$  & $1.8\%$ & $6\%$ \\
$8$ & $512$  & $0.8503$ & $0.2694$ & $2.26$ & $0.04\%$ & $2.9\%$ & $25\%$ \\
$8$ & $2048$ & $0.8507$ & $0.2793$ & $2.72$ & $0.08\%$ & $0.6\%$ & $9\%$ \\
\bottomrule
\end{tabular}
\end{minipage}
\end{center}

\begin{figure}[t]\centering
\includegraphics[width=\linewidth]{fig_em_innovation.pdf}
\caption[What EM recovers from noised sequences]{\emph{Left:} fitted $C=8$
mixture innovation densities for three data budgets against the Laplace
truth; the discrepancies sit at the cusp and in the tails. \emph{Right:}
relative errors in $\alpha$, innovation variance and fourth cumulant against
$\nseq$, the hierarchy from correlation to scale to shape.}
\label{fig:em-innovation}
\end{figure}
